\PassOptionsToPackage{svgnames}{xcolor}
\documentclass[12pt]{article}



\usepackage[margin=1in]{geometry}  
\usepackage{graphicx}             
\usepackage{amsmath}              
\usepackage{amsfonts}              
\usepackage{framed}               
\usepackage{amssymb}
\usepackage{array}
\usepackage{amsthm}
\usepackage[nottoc]{tocbibind}
\usepackage{bm}
\usepackage[object=vectorian]{pgfornament} 
\usepackage{enumitem}

\colorlet{shadecolor}{lightgray!25}
\newcommand{\sectionline}{%
  \noindent
  \begin{center}
  {\color{DarkViolet}
    \resizebox{0.5\linewidth}{1ex}
    {{%
    {\begin{tikzpicture}
    \node  (C) at (0,0) {};
    \node (D) at (9,0) {};
    \path (C) to [ornament=85] (D);
    \end{tikzpicture}}}}}%
    \end{center}
  }

  \newcommand\norm[1]{\left\lVert#1\right\rVert}
\setlength{\parindent}{0cm}
\setlength{\parskip}{0em}
\newcommand{\Lim}[1]{\raisebox{0.5ex}{\scalebox{0.8}{$\displaystyle \lim_{#1}\;$}}}
\newtheorem{definition}{Definition}[section]
\newtheorem{theorem}{Theorem}[section]
\newtheorem{notation}{Notation}[section]
\theoremstyle{definition}
\newcommand\fl[1]{\text{fl}\left(#1\right)}
\DeclareMathOperator{\arcsec}{arcsec}
\DeclareMathOperator{\arccot}{arccot}
\DeclareMathOperator{\arccsc}{arccsc}
\DeclareMathOperator{\spn}{Span}
\DeclareMathOperator{\R}{\mathbf{R}}
\DeclareMathOperator{\diff}{d}
\setcounter{tocdepth}{1}
\begin{document}

\title{Revision notes - MA2213}
\author{Ma Hongqiang}
\maketitle
\tableofcontents

\clearpage
%\twocolumn
\section{Computer Arithmetic and Computational Errors}
\subsection{Number Systems}
\begin{definition}[Decimal System]
\hfill\\\normalfont The number system we are used to is the \textbf{decimal system} and the number "10" which plays an important role is called the \textbf{base} of the decimal system.
\end{definition}
In general, we can take other positive integer $N>1$ as a base. In particular, $N = 2, 8, 16$ are used in most digital computers and the systems with $N = 2,8,16$ are known as \textbf{binary}, \textbf{octal}, \textbf{hexadecimal} number systems respectively.\\
\subsubsection{Decimal to Binary}
The conversion of decimal number to binary is performed in two steps.\\
Take $(53.7)_{10}$ as an example.
\begin{enumerate}
  \item \textbf{Integer Part}. Convert decimal integers to binary by deviding by 2 successively and recording the remainders. \\The remainders, 0 or 1, are recorded by starting at the decimal point and moving away to the left. For$ (53)_{10}$, we have
  \begin{alignat*}{3}
    53/2 &= 26 &&\R 1\\
    26/2 &= 13 &&\R 0\\
    13/2 &= 6 &&\R 1\\
    6/2 &=3&&\R 0\\
    3/2 &= 1&&\R 1\\
    1/2 &= 0&&\R 1
  \end{alignat*}
  Therefore, the base 10 number 53 can be written as binary
  \[
(110101)_2 = 2^5+2^4+2^2+2^0 = 53
  \]
  \item \textbf{Fractional Part}. Convert $(0.7)_{10}$ to binary as follows: \\Multiply by 2 successfully and record the integer parts, moving away from the decimal point to the right:
  \begin{alignat*}{3}
    .7\times2 &= .4 &&+ 1\\
    .4\times2 &= .8 &&+ 0\\
    .8\times2 &= .6 &&+ 1\\
    .6\times2 &=.2&&+ 1\\
    .2\times2 &= .4&&+ 0\\
    .4\times2 &= .8&&+ 0\\
    \cdots&\cdots\cdots&& 
  \end{alignat*}
  Notice that the process repeats after four steps and will repeat infinitely exactly the same way. Therefore,
  \[
(0.7)_{10} = (.1\overline{0110})_2
  \]
  \item Hence,
  \[
(53.7)_{10} = (110101.1\overline{0110})_2
  \]
\end{enumerate}
\subsubsection{Binary to Decimal}
The conversion of binary number to decimal needs to tackle nonterminating binary numbers.\\
Take $x=(.10\overline{101})_2$ as an example.\\
Multiplying by $2^2$ shifts $x$ to
\[
y:=2^2 x = (10.\overline{101})_2 = (10)_2+(.\overline{101})_2 = (2)_{10}+(.\overline{101})_2
\]
The fractional part of $y$, i.e.,
\[
z:=(.\overline{101})_2
\]
is calculated as follows:
\[
2^3z = (101.\overline{101})_2
\]
So,
\[
(2^3-1)z= (101)_2 = (5)_{10}
\]
i.e.
\[
z = \left(\frac{5}{7}\right)_{10}
\]
Thus, 
\[
y=(2)_{10}+\left(\frac{5}{7}\right)_{10}=\left(\frac{19}{7}\right)_{10}
\]
and hence,
\[
x=\frac{y}{2^2}=\left(\frac{19}{28}\right)_{10}
\]
\begin{theorem}[Representation of Numbers]
\hfill\\\normalfont In general, any \textbf{natural numbers} $N\geq 2$ can be used as base.\\
Every positive real number $a$ has a unique representation of the form
\[
a=a_mN^m+a_{m-1}N^{m-1}+\cdots+a_1N^1+a_0N^0+a_{-1}N^{-1}+a_{-2}N^{-2}+\cdots
\]
where
\[
0\leq a_i\leq N-1,\;\;\;a_m\neq 0
\]
Equivalently,
\[
a=\sum_{i=0}^\infty a_{m-i}N^{m-i}
\]
where
\[
0<a_m<N; 0\leq a_{m-i}\leq N - 1\;\;\; \forall i\geq 1
\]
The above non-terminating representation of number $a$ is not feasible in practical computation. We onl deal with, in practice, numbers which have \textit{terminating} representation, i.e. of the following form
\[
\alpha = \sum_{i = 0}^{n-1}\alpha_{m-i}N^{m-i}
\]
where $0<\alpha_m<N$ and $0\leq \alpha_{m-i}<N$ for $i = 1,2,\ldots, n-1$.
\end{theorem}
\begin{definition}[Significant decimal digits]
\hfill\\\normalfont If $N = 10$, then the numbers $\alpha_{m-i}, 0\leq i \leq n-1$ are called \textbf{significant decimal digits} and the number $\alpha$ is said to have $n$ significant decimal digits.
\end{definition}
\subsection{Chopping and Rounding}
If the number $a$ has more significant digits than what we wish to have, then we need to replace $a$ by an approximate number $a^\ast$ which contains a smaller number of digits, say $n$. There are two ways of terminating the number $a$ to a given significant number of digits, namely \textbf{chopping} or \textbf{rounding}.
\begin{definition}[Chopping]
\hfill\\\normalfont In \textbf{chopping}, we retain only the first $n$ significant digits in the number $a$.
\end{definition}
\begin{definition}[Rounding]
\hfill\\\normalfont In \textbf{rounding}, the following rules are usually practiced:
\begin{itemize}
  \item Retain the first $n$ significant digits, and
  \item if the $(n+1)$th significant digit is less than $5$, leave the $n$th significant digit unchanged;
  \item otherwise, if the $(n+1)$th significant digit is greater or equal to $5$, add unity to the $n$th significant digit.
\end{itemize}
\end{definition}
\subsection{Floating Point Representation}
\begin{definition}[Floating Point Numbers]
\hfill\\\normalfont Number of the form
\[
\pm(0.a_1a_2a_3\ldots a_m)\times N^e, \text{with }a_1\neq 0
\]
are called \textbf{floating point numbers}. The factor $0.a_1a_2a_3\ldots a_n$ is called the \textbf{mantissa}, $e$ is called the \textbf{exponent} and $N$ is the \textbf{base}.
\end{definition}
The folating-point mode is for storing real numbers.
\subsection{Basic Concepts in Error Estimation}
\begin{definition}[Sources of Error]
\hfill\\\normalfont Numerical results are affected by many types of error.
\begin{itemize}
  \item \textbf{Error in Given Input Data}\\The input data can be \textit{result of measurements} which are inexact, or produced by some \textit{arithmetic process using round-off process}.
  \item \textbf{Round-off Errors During Computation}\\This is due to working with \textit{finite machine precision}.
  \item \textbf{Truncation Error}\\Consider the Taylor series xpansion for $\sin(x)$ about $x=0$
  \[
\sin(x) = x-\frac{x^3}{3!}+\frac{x^5}{5!}-\frac{x^7}{7!}+\ldots
  \]
which is an infinite series.\\Truncation error occurs when the \textit{infinite series is broken off after a finite number of terms}.\\In general, truncation errors are error committed when \textit{a limiting process is truncated before one has come to the limiting value}.
  \item Simplifications in the Mathematical Model
  \item Human Error
\end{itemize}
\end{definition}
\subsection{Absolute and Relative Error}
\begin{definition}[Absolute Error]
\hfill\\\normalfont Let $x^\ast$ be an approximation to $x$, then the \textbf{absolute error} in approximating $x$ is given by $|x-x^\ast|$.
\end{definition}
From the above definition, the floating point representation of $x$, $\fl{x}$ has absolute error
\[
|x-\fl{x}|
\]
\begin{definition}[Absolute Error Bound]
\hfill\\\normalfont Any non-negative number $\delta(x^\ast)$ satisfying the inequality
\[
|x-x^\ast|<\delta(x^\ast)
\]
is called an \textbf{absolute error bound}.
\end{definition}
\begin{definition}[Decimal Places]
\hfill\\\normalfont The two numbers $x$ is said to agree to $k$ \textbf{decimal places} if $k$ is the largest non-negative integer such that
\[
|x-x^\ast|\leq 0.5\times 10^{-k}
\]
\end{definition}
\begin{definition}[Relative Error]
\hfill\\\normalfont If $x^\ast$ is an approximation to $x$, then the \textbf{relative error} is defined by
\[
\frac{|x-x^\ast|}{|x|}\;\;\;\text{provided that } x\neq 0
\]
\end{definition}
For a $t$-digit machine,
\[
\frac{|x-\fl{x}|}{|x|}\leq \frac{1}{2}\times 10^{1-t}
\]
in rounding mode and
\[
\frac{|x-\fl{x}|}{|x|}\leq 10^{1-t}
\]
in chopping mode.
\subsection{Computation in Floating Point Arithmetic}
Take the example of addition/subtraction of two floating point number $0.a_1a_2a_3\ldots a_m\times 10^{N_1}\pm 0.b_1b_2b_3\ldots b_m\times 10^{N_2}$.
\begin{align*}
&0.a_1a_2a_3\ldots a_m\times 10^{N_1}\pm 0.b_1b_2b_3\ldots b_m\times 10^{N_2}\\
=&\begin{cases}
(0.a_1a_2a_3\ldots a_m\pm 0.b_1b_2b_3\ldots b_m)\times 10^{N_1}&\text{if }N_1= N_2\\
(0.a_1a_2a_3\ldots a_m 0\ldots 0\pm 0.0\ldots 0b_1b_2b_3\ldots b_m)\times 10^{N_1}&\text{if }N_1>N_2
\end{cases}
\end{align*}
\subsection{Propagation of Errors in Function Evaluation}
\subsubsection{Function of one variable}
In finding the value of $f(x)$ by approximating $f(x^\ast)$, where $x^\ast$ is the known approximating value of $x$,\\let $\delta(x^\ast)$ be an absolute error bound for $|x-x^\ast|$, i.e.,
\[
|x-x^\ast|\leq \delta(x^\ast)
\]
If the function $f(x)$ is differentiable, then by the mean-value theorem, we have
\[
f(x)-f(x^\ast)=f^\prime(p)(x-x^\ast)
\]
where $p$ is some number between $x$ and $x^\ast$.\\
Therefore, 
\[
\delta(f(x^\ast))\leq \max_{t\in\mathbf{I}}|f^\ast(t)|\delta(x^\ast)
\]
where $\mathbf{I}$ is the interval $(x,x^\ast)$ if $x<x^\ast$ or $(x^\ast,x)$.\\
Sometimes, upper bound of $f^\prime$ is unavailable. Yet, with the following assumption:
\begin{itemize}
  \item $\delta(x^\ast)$ is small
  \item $f^\prime(x^\ast)\neq 0$, and
  \item $f^\prime(t)$ is nearly constant  near $x^\ast$ (i.e., $f^\prime(t)$ does not very greatly for $t$ between $x$ and $x^\ast$)
\end{itemize}
We have $D\approx |f^\prime(x^\ast)|$, and consequently
\[
\delta(f(x^\ast))\approx|f^\prime(x^\ast)|\delta(x^\ast)
\]
\subsubsection{Function of Several Variables}
\begin{theorem}[Error bounds of function of several variables]
\hfill\\\normalfont
Generalisation suggests, if $x_i^\ast$ is an estimate of $x_i$, $1\leq i\leq n$, then with $x^\ast=(x_1^\ast,x_2^\ast,\ldots,x_n^\ast)$, we have
\[
\delta(f(x^\ast=(x_1^\ast,x_2^\ast,\ldots,x_n^\ast))\approx\left\lvert\frac{\partial f(x^\ast)}{\partial x_1}\right\rvert \delta(x_1^\ast)+\left\lvert\frac{\partial f(x^\ast)}{\partial x_2}\right\rvert \delta(x_2^\ast)+\cdots+\left\lvert\frac{\partial f(x^\ast)}{\partial x_n}\right\rvert \delta(x_n^\ast)
\]
\end{theorem}
\subsection{Catastrophic Cancellation}
Calculations involving the \textbf{subtraction} of two \textit{nearly equal} numbers can result in considerable loss of accuracy due to cancellation.\\
Let $x_1$ and $x_2$ be two nearly equal numbers and denote the error in $x_1$ and $x_2$ by $\delta(x_1)$ and $\delta(x_2)$, respectively, we have
\begin{alignat*}{3}
&&y&=x_1-x_2\\
\Rightarrow&&|\delta(y)|&\leq |\delta(x_1)|+|\delta(x_2)|\\
\Rightarrow&&\left\lvert\frac{\delta(y)}{y}\right\lvert&\leq \frac{|\delta(x_1)|+|\delta(x_2)|}{|x_1-x_2|}
\end{alignat*}
Specifically, suppose two nearly equal numbers $x_1$ and $x_2$, having the same exponent $n$ and being in the $k$ digi representations, are represented in floating point form:
\begin{align*}
x_1&=0.a_1a_2\cdots a_p\alpha_{p+1}\cdots\alpha_{k}\times 10^n\\
x_2&=0.a_1a_2\cdots a_p\beta_{p+1}\cdots\beta_{k}\times 10^n
\end{align*}
Then their difference $x_1-x_2$ in floating point representationwill have some  digits in the mantissa to be $0$, i.e.,
\[
x_1 - x_2 = \fl{\fl{x_1}-\fl{x_2}}=0.\gamma_{p+1}\cdots\gamma_{k}\underbrace{0\cdots 0}_{p\;0\text{'s}}\times 10^{n-p}
\]
So, in floating point calculation, $x_1-x_2$ has at most $k-p$ significant digits; $p$ significant digits have been lost or canceled due to the subtraction.
\subsubsection{Ways to Avoid Subtraction} 
\begin{theorem}[$\sqrt{x+\varepsilon}-\sqrt{x}$]
\hfill\\\normalfont If $|\varepsilon|\ll x$, then 
\[
\sqrt{x+\varepsilon}-\sqrt{x}=\frac{\varepsilon}{\sqrt{x+\varepsilon}+\sqrt{x}}
\]
presents a better way of calculation that reduces loss of accuracy.
\end{theorem}
\begin{theorem}[Roots of $ax^2+bx+c=0$]
\hfill\\\normalfont Consider the quadratic equation
\[
ax^2+bx+c = 0\;\;\;a\neq 0
\]
The roots of equations are
\[
x_1=\frac{-b+\sqrt{b^2-4ac}}{2a}\;\;\;\text{and}\;\;\;x_2=\frac{-b-\sqrt{b^2-4ac}}{2a}
\]
To reduce loss of accuracy, some alternative pathway include:
\begin{itemize}
\item $x_1=\frac{-2c}{b+\sqrt{b^2-4ac}}\;\;\;\text{and}\;\;\;x_2=\frac{2c}{-b+\sqrt{b^2-4ac}}$
\item $x_1+x_2 = -b$
\item $x_1x_2=c$
\end{itemize}
\end{theorem}
\begin{theorem}[Evaluation of polynomial $f(x)$]
\hfill\\\normalfont First note that $ax^n = \fl{\fl{a}\times\fl{x^n}}$ where $\fl{x^n}=\fl{\fl{x^{n-1}}\times\fl{x}}$ recursively.\\Instead of calculating $f(x) = \fl{\sum_{i=0}^n a_ix^i}$, nesting may provide better result:
\[
f(x)=\fl{(\cdots(a_nx+a_{n-1})x+\cdots)x+a_0}
\]
\end{theorem}
\subsection{Numerical Instability}
\begin{definition}[Numerical Stability]
\hfill\\\normalfont An algorithm is said to be \textbf{stable} if the effect of its errors on the final result is \textbf{negligible}.
\end{definition}
\clearpage
\section{Numerical Solution of Linear Systems of Equations}
Definitions and method of computation of matrix can be found in \texttt{MA2101.pdf}, and are therefore omitted here.
\subsection{System of Linear Equations}
In this chapter, only linear systems with \textit{unique} solution is concerned.\\
Consider the linear system
\[
Ax=b
\]
where $A\in\mathbb{M}_{n}(\mathbb{R})$ and $b\in\mathbb{R}_c^n$. If $A$ is invertible, a \textit{unique} solution $x$ exists and given by
\[
x=A^{-1}b
\]
There are two classes of method of solving $Ax=b$ numerically:
\begin{itemize}
  \item \textbf{Direct} methods
  \item \textbf{Iterative} methods
\end{itemize}
In this chapter, only direct methods are studied.\\
Two basic direct method include
\begin{itemize}
  \item Cramer's rule
  \item Gaussian Elimination
\end{itemize}
\subsubsection{Direct Method: Cramer's Rule}
\begin{definition}[Cramer's Rule]
\hfill\\\normalfont By \textbf{Cramer's rule}, the $n\times n$ linear system
\[
Ax=b
\]
has solution
\[
x_i=\frac{d_i}{d} \;\;\; i = 1,\ldots, n
\]
where $d:=\det(A)\neq 0, d_i=\det(A_i)$, $A_i$ the matrix obtained by replacing the $i$-th column of $A$ by $b$.\\
The determinant $d$ is given by \textbf{Laplace Theorem}:
\[
d=(-1)^{i+1}a_{i1}D_{i1}+(-1)^{i+2}a_{i2}D_{i2}+\cdots+(-1)^{i+n}a_{in}D_{in}
\]
where $D_{ij}$ is the determinant of the submatrix obtained from $A$ by deleting its $i$th row and $j$th column.
\end{definition}
\begin{theorem}[Computational Complexity of Cramer's Rule]
\hfill\\\normalfont Suppose it needs $m_n$ operations of multiplication to compute the determinant of a $n\times n$ matrix. Then, we have
\[
m_n=n+nm_{n-1}\;\;\;m_1=1
\]
Here, the first term $n$ is results from multiplication of $a_{ij}$ and $D_{ij}$.\\
Thus,  
\begin{align*}
m_n &=n+nm_{n-1}=n+n[(n-1)+(n-1)m_{n-2}]\\
&=n+n(n-1)+n(n-1)(n-2)+\cdots+n(n-1)\cdots 3\cdot2\\
&>n! 
\end{align*}
Hence, in order to solve an $n\times n$ linear system by Cramer's rule and Laplace Theorem, we have do at least
\[
(n+1)n!=(n+1)!
\]
  multiplication.\\Clearly, Cramer's rule is too \textit{computationally expensive}.
\end{theorem}
Also, solving nonlinear linear system by \textbf{matrix inversion} is \textit{computationally expensive} and often leads to more \textit{inaccuracies}.
\subsection{Gaussian elimination}
\begin{definition}[Elementary Row Operations]
\hfill\\\normalfont Let $E_i$ be the $i$th equation in a linear system $Ax=b$. The three \textbf{elementary row operations} permitted to solve this linear system are:
\begin{itemize}
\item $E_i\leftarrow \lambda E_i$, where $\lambda$ is a nonzero constant.
\[
\begin{bmatrix}
1&&&&\\
&\ddots&&&\\
&&\lambda&&\\
&&&\ddots&\\
&&&&1
\end{bmatrix}\begin{bmatrix}
a_{11}&\cdots&a_{1n}\\
\vdots&&\vdots\\
a_{i1}&\cdots&a_{in}\\
\vdots&&\vdots\\
a_{n1}&\cdots&a_{nn}
\end{bmatrix}=\begin{bmatrix}
a_{11}&\cdots&a_{1n}\\
\vdots&&\vdots\\
\lambda a_{i1}&\cdots&\lambda a_{in}\\
\vdots&&\vdots\\
a_{n1}&\cdots&a_{nn}
\end{bmatrix}
\]
\item $E_i\leftarrow E_i-lE_j$, where $l$ is a nonzero constant.
\[
\begin{bmatrix}
1&&&&\\
&\ddots&&&\\
&-l&1&&\\
&&&\ddots&\\
&&&&1
\end{bmatrix}^{\footnotemark}\begin{bmatrix}
a_{11}&\cdots&a_{1n}\\
\vdots&&\vdots\\
a_{i1}&\cdots&a_{in}\\
\vdots&&\vdots\\
a_{n1}&\cdots&a_{nn}
\end{bmatrix}=\begin{bmatrix}
a_{11}&\cdots&a_{1n}\\
\vdots&&\vdots\\
 a_{i1}-la_{j1}&\cdots&a_{in}-la_{jn}\\
\vdots&&\vdots\\
a_{n1}&\cdots&a_{nn}
\end{bmatrix}
\]\footnotetext{$-l$ is at $i$th row and $j$th column of the leftmost matrix.}
\item $E_i\leftrightarrow E_j$
\[
\begin{bmatrix}
1&&&&\\
&\ddots&&1&\\
&&1&&\\
&1&&\ddots&\\
&&&&1
\end{bmatrix}\begin{bmatrix}
a_{11}&\cdots&a_{1n}\\
a_{i1}&\cdots&a_{in}\\
\vdots&\cdots&\vdots\\
a_{j1}&\cdots&a_{jn}\\
a_{n1}&\cdots&a_{nn}
\end{bmatrix}=\begin{bmatrix}
a_{11}&\cdots&a_{1n}\\
a_{j1}&\cdots&a_{jn}\\
\vdots&\cdots&\vdots\\
a_{i1}&\cdots&a_{in}\\
a_{n1}&\cdots&a_{nn}
\end{bmatrix}
\]
\end{itemize}
\end{definition}
By a sequence of the above operations, a linear system can be transformed to a more easily solved linear system with the \textbf{same} set of solutions.\\
The Gaussian elimination for solving $Ax=b$ involves 2 stages/
\begin{itemize}
\item A forward course where a sequence of elementary row operations are applied to $\mathbf{A}$ to reduce it to an upper triangular form $\mathbf{U}$.
\item Backward substitution process
\end{itemize}
\begin{theorem}[General method of Gaussian Elimination on $\mathbf{A}\mathbf{x}=\mathbf{B}$]
\hfill\\\normalfont Let the original system be denoted by 
\[
\mathbf{A}^{(0)}\mathbf{x}=\mathbf{b}^{(0)}
\]
where $\mathbf{A}^{(0)} = \mathbf{A} = (a_{ij}^{(0)})$, $\mathbf{b}^{(0)}=\mathbf{b} = \begin{bmatrix}
b_1^{(0)}\\
\vdots\\
b_i^{(0)}\\
\vdots\\
b_n^{(0)}
\end{bmatrix}$\\
Step 1: If $a_{11}^{(0)}\neq 0$, let
\begin{alignat*}{23}
l_{i,1}&=\frac{a^{(0)}_{i,1}}{a^{(0)}_{1,1}}\footnotemark&&\;\;\;\;\;i = 2,3,\ldots, n\\
a^{(1)}_{i,j}&=a^{(0)}_{i,j}-l_{i,1}\times a^{(0)}_{1,j}\footnotemark&&\;\;\;\;\;i,j = 2,3,\ldots,n\\
b^{(1)}_{i}&=b^{(0)}_i-l_{i,1}\times b^{(0)}_1&&\;\;\;\;\;i = 2,\ldots,n
\end{alignat*}
\footnotetext{In general, $l_{p,q}$ is used in the operations on $q$th column, and denotes the ratio between the entry on the $p$th row($p>q$),$a^{(q-1)}_{p,q}$, to the $q$th diagonal entry, $a^{(q-1)}_{q,q}$.}
\footnotetext{In general, $a_{i,j}^{(k)}$,$(i<j)$ is obtained by operations on the $k$th column, and is the result of zeroing the $a^{(k)}_{i,k}$ entry($i>k$), so it equals the previous entry $a_{i,j}^{(k-1)}$ minus $l_{i,k}$ times the entry on the $k$th row $a^{(k-1)}_{k,j}$.}
Then we obtain the equivalent system
\[
\mathbf{A}^{(1)}\mathbf{x}=\mathbf{b}^{(1)}
\]
where
\[
\mathbf{A}^{(1)}=\begin{bmatrix}
a_{11}^{(0)}&a_{12}^{(0)}&\cdots&a_{1n}^{(0)}\\
0&a_{22}^{(1)}&\cdots&a_{2n}^{(1)}\\
0&\vdots&\ddots&\vdots\\
0&a_{n2}^{(1)}&\cdots&a_{nn}^{(1)}\end{bmatrix},\;\;\; \mathbf{b}^{(1)} = \begin{bmatrix}
b_1^{(0)}\\
b_2^{(1)}\\
\vdots\\
b_n^{(1)}
\end{bmatrix}
\]
Step 2: If $a_{22}^{(1)}\neq 0$, repeat Step 1 to eliminate $\mathbf{x}_2$ from row 3 to $n$.\\After $k$ steps, we obtain the equivalent system
\[
\mathbf{A}^{(k)}\mathbf{x}=\mathbf{b}^{(k)}
\]
where $\mathbf{A}^{(k)}$ takes the form
\[
\mathbf{A}^{(k)}=\begin{bmatrix}
a_{11}^{(0)}&\cdots       &\cdots&\cdots            &\cdots&a_{1n}^{(0)}\\
            &a_{22}^{(1)} &\cdots&\cdots            &\cdots&a_{2n}^{(1)}\\
            &             &\ddots&\vdots            &\cdots&\cdots      \\
            &             &      &a_{k+1,k+1}^{(k)} &\cdots&a_{k+1,n}^{(k)}\\
            &             &      &\vdots            &\ddots&\vdots\\
            &             &      &a_{n,k+1}^{(k)} &\cdots&a_{n,n}^{(k)}
            \end{bmatrix}
\]
If $a_{k+1,k+1}^{(k)}\neq 0$, then let
\begin{align*}
l_{i,k+1}&=\frac{a^{(k)}_{i,k+1}}{a^{(k)}_{k+1,k+1}}\\
a^{(k+1)}_{i,j}&=a^{(k)}_{i,j}-l_{i,k+1}\times a^{(k)}_{k+1,j}\\
b_i^{(k+1)}&=b_i^{(k)}-l_{i,k+1}b^{(k)}_{k+1}
\end{align*}
Finally, after at most $n-1$ steps, the system becomes
\[
\mathbf{A}^{(n-1)}\mathbf{x}=\mathbf{b}^{(n-1)}
\]
which is an upper triangular system, where
\[
\mathbf{U}=\mathbf{A}^{(n-1)}=\begin{bmatrix}
a_{11}^{(0)}&a_{12}^{(0)}&\cdots&a_{1n}^{(0)}\\
 &a_{22}^{(1)}&\cdots&a_{2n}^{(1)}\\
 & &\ddots&\vdots\\
&&&a_{nn}^{(n-1)}\end{bmatrix}
\]and\[
\mathbf{b}^{(n-1)} = \begin{bmatrix}
b_1^{(0)}\\
b_2^{(1)}\\
\vdots\\
b_n^{(n)}
\end{bmatrix}
\]
Step 3: The solution vector $\mathbf{x}$ is solved by back substitution:
\[
x_n = \frac{b_n^{(n-1)}}{a_{nn}^{(n-1)}}
\]
and for $k = n-1, n-2,\ldots, 1$,
\[
x_k = \frac{1}{a_{kk}^{(k-1)}}\left(b_k^{(k-1)}-\sum_{j=k+1}^n a_{kj}^{(k-1)}x_j\right)
\]
The above algorithm is known as \textbf{Gaussian elimination}.
\end{theorem}
\subsection{Triangular Factorisation: $\mathbf{A}=\mathbf{LU}$}
We note, in the $(k+1)$th step of above Gaussian elimination, all rows below the $(k+1)$th row are minused by a multiple of the $(k+1)$th row. This operation can be viewed by a collection of elementary row operation 2 and represented by
\[
\mathbf{A}^{(k+1)}=\mathbf{L}^{(k)}\mathbf{A}^{(k)}
\]
where
\[
\mathbf{L}^{(k)}=\begin{bmatrix}
1& & & & &\\
 &\ddots&&&&\\
 &&1&&&\\
&&-l_{k+2,k+1}&\ddots&&\\
&&\vdots&&\ddots&\\
&&-l_{n,k+1}&&&1
\end{bmatrix}
\]
This relation gives
\[
\mathbf{U}=\mathbf{A}^{(n-1)}=\mathbf{L}^{(n-2)}\mathbf{L}^{(n-3)}\cdots\mathbf{L}^{(0)}\mathbf{A}^{(0)}
\]
Also, note the inverse of $\mathbf{L}^{(k)}$ is
\[
(\mathbf{L}^{(k)})^{-1}=\begin{bmatrix}
1& & & & &\\
 &\ddots&&&&\\
 &&1&&&\\
&&l_{k+2,k+1}&\ddots&&\\
&&\vdots&&\ddots&\\
&&l_{n,k+1}&&&1
\end{bmatrix}
\]
a flip of signs of all $l_{i,k+1}$ from $\mathbf{L}^{(k)}$.\\
Therefore,
\[
(\mathbf{L}^{(0)})^{-1}(\mathbf{L}^{(1)})^{-1}\cdots (\mathbf{L}^{(n-2)})^{-1} \mathbf{U}=\mathbf{A}
\]
Note that $(\mathbf{L}^{(0)})^{-1}(\mathbf{L}^{(1)})^{-1}\cdots (\mathbf{L}^{(n-2)})^{-1}$ is upper triangular and is defined as $\mathbf{L}$:
\[
(\mathbf{L}^{(0)})^{-1}(\mathbf{L}^{(1)})^{-1}\cdots (\mathbf{L}^{(n-2)})^{-1}=\begin{bmatrix}
1&&&&&\\
l_{21}&\ddots&&&&\\
\vdots&&1&&&\\
l_{k+2,1}&\cdots&l_{k+2,k+1}&\ddots&&\\
\vdots&&\vdots&&\ddots&\\
l_{n,1}&\cdots&l_{n,k+1}&\cdots&l_{n,n-1}&1
\end{bmatrix}:=\mathbf{L}
\]
Hence, $\mathbf{A}=\mathbf{L}\mathbf{U}$, where $\mathbf{U}=\mathbf{A}^{(n-1)}$ and $\mathbf{L}$ is a lower triangular matrix.
\begin{theorem}\normalfont If $\mathbf{U}$ and $\mathbf{L}$ are the upper and lower triangular matrices defined above, then $\mathbf{A}=\mathbf{L}\mathbf{U}$.
\end{theorem}
\subsection{Compact Forms of Gaussian Elimination}
If $\mathbf{A}$ admits an $\mathbf{LU}$ decomposition, we can solve the system of equations $\mathbf{LUx}=\mathbf{b}$ in two stages.
\begin{itemize}
  \item Set $\mathbf{z}:=\mathbf{Ux}$. Solve $\mathbf{Lz}=\mathbf{b}$ for $\mathbf{z}$.
  \item Solve $\mathbf{Ux}=\mathbf{z}$ for $\mathbf{x}$
\end{itemize}
\begin{theorem}[Existence and Uniqueness of $\mathbf{LU}$ decomposition]
\hfill\\\normalfont Let $\mathbf{A}$ be an $n\times n$ matrix and $\mathbf{A}^{(k)}$ be the $k\times k$ matrix formed from the first $k$ rows and columns of $\mathbf{A}$. If $\det(\mathbf{A}^{(k)})\neq 0 \forall k = 1,2,\ldots, n-1$, then there exists a unique lower triangular matrix $\mathbf{L}=(l_{ij})$ with $l_{ii}=1 \forall i=1,2\ldots,n$ and a unique upper triangular matrix $\mathbf{U}=(u_{ij})$ such that $\mathbf{A}=\mathbf{LU}$.
\end{theorem}
\subsection{Compact Forms of Gaussian Elimination}
The $\mathbf{LU}$ decomposition can be calculated directly by Gaussian Elimination.\\
Given that $mathbf{A}$ admits a $\mathbf{LU}$ decomposition $\mathbf{A}=\mathbf{LU}$, we have, entry-wise
\[
a_{ij} = \sum_{k=1}^n l_{ik}u_{kj}\;\;\;i,j = 1,2,\ldots,n
\]
Note that the above system has $n^2$ equations and $n^2+n$ unknowns $l_{ik}, i\geq k$ and $u_{kj},k\leq j$. Thus $n$ unknowsn may be set arbitrarily.
\begin{theorem}[Doolittle Method]
\hfill\\\normalfont For Doolittle Method, set diagonal entries of $\mathbf{L}$, $l_{kk}:=1, k = 1,2,\ldots n$ and assume that $u_{kk}\neq 0,\forall k \in[1,n]\cap\mathbb{Z}$.\\
The sequence of calculation is as follows:
\begin{enumerate}
  \item First row of $\mathbf{U}$: $u_{1k}, 1\leq k\leq n$
  \item First column of $\mathbf{L}$: $l_{k1}, 2\leq k\leq n$\footnote{if $k = 1, l_{k1}=l_{11}=1$ by definition}
  \item Second row of $\mathbf{U}$: $u_{2k}, 2\leq k\leq n$.
  \item[] $\vdots$
\end{enumerate}
Obviously, when calculating $k$th row of $\mathbf{U}$, row 1 to $k-1$ of $\mathbf{U}$ and column 1 to $k-1$ of $\mathbf{L}$ will be known.\\
Also, when calculating $k$th row of $\mathbf{L}$, row 1 to $k$ of $\mathbf{U}$ and column 1 to $k-1$ of $\mathbf{L}$ will be known.\\
Suppose $u_{kj}, (j\geq k)$ is of concern, we have
\[
a_{kj}=\sum_{r=1}^k l_{kr}u_{rj}
\]
Substitute in $l_{kk}=1$ and rearranging,
\[
u_{kj} = a_{kj} -\sum_{r=1}^{k-1}l_{kr}u_{rj}
\]
Afterwards, suppose $l_{ik}, (i > k)$ is of concern, we have
\[
a_{ik} = \sum_{r=1}^k l_{ir}u_{rk}
\]
Rearranging,
\[
a_{ik}=\sum_{r=1}^{k-1} l_{ir}u_{rk}+l_{ik}u_{kk}
\]
Therefore,
\[
l_{ik}=\frac{a_{ik}-\sum_{r=1}^{k-1} l_{ir}u_{rk}}{u_{kk}}
\]
The above two equation can be interleaved to obtain $\mathbf{L}$ and $\mathbf{U}$.
\end{theorem}
\begin{theorem}[Crout Method]
\hfill\\\normalfont For Crout Method, set diagonal entries of $\mathbf{U}$, $u_{ii}=1, k = 1,2,\ldots, n$ and assume that $l_{kk}\neq 0 \forall k\in[1,n]\cap\mathbb{Z}$.\\
The sequence of calculation is as follows:
\begin{enumerate}
  \item First column of $\mathbf{L}$: $l_{k1}, 1\leq k\leq n$
  \item First row of $\mathbf{U}$: $u_{1k}, 2\leq k\leq n$
  \item Second column of $\mathbf{L}$: $l_{k2}, 2\leq k\leq n$
  \item[] $\vdots$
\end{enumerate}
Obviously, when calculating $k$th column of $L$, column 1 to $k-1$ of $\mathbf{L}$ and row 1 to $k-1$ of $\mathbf{U}$ will be known.\\
Also, when calculating $k$th row of $\mathbf{U}$, column 1 to $k$ of $\mathbf{L}$ and row 1 to $k-1$ of $\mathbf{U}$ will be known.\\
Suppose $l_{ik}, (i\geq k)$ is of concern, we have
\[
a_{ik} = \sum_{r = 1}^k l_{ir}u_{rk}
\]
Substitute in $u_{kk}=1$ and rearranging,
\[
l_{ik} = a_{ik} - \sum_{r = 1}^{k-1}l_{ir}u_{rk}
\]
Afterwards, suppose $u_{kj}, (j\geq k+1)$ is of concern, we have
\[
a_{kj} = \sum_{r = 1}^k l_{kr}u_{rj}
\]
Rearranging, 
\[
u_{kj} = \frac{a_{kj}-\sum_{r=1}^{k-1}l_{kr}u_{rj}}{l_{kk}}
\]
The above two equation can be interleaved to obtain $\mathbf{L}$ and $\mathbf{U}$.
\end{theorem}
\subsection{$\mathbf{LU}$ Decomposition for Tridiagonal Matrices}
\begin{theorem}\hfill\\\normalfont
Let $\mathbf{A}$ be a \textbf{tridiagonal matrices} 
\[
\mathbf{A}=\begin{bmatrix}
a_1&c_1&&&&\\
b_2&a_2&c_2&&&\\
&\ddots&\ddots&\ddots&&\\
&&\ddots&\ddots&\ddots&\\
&&&b_{n-1}&a_{n-1}&c_{n-1}\\
&&&&b_n&a_n
\end{bmatrix}
\]
If an $\mathbf{LU}$ decomposition exists for $\mathbf{A}$, then
\[
\mathbf{A}=\mathbf{LU}
\]
where
\[
\mathbf{L}=\begin{bmatrix}
1&&&&&\\
\beta_2&1&&&&\\
&\ddots&\ddots&\ddots&&\\
&&\ddots&\ddots&\ddots&\\
&&&\beta_{n-1}&1&\\
&&&&\beta_n&1
\end{bmatrix}\;\;\;\text{and}\;\;\;\mathbf{U}=\begin{bmatrix}
\alpha_1&c_1&&&&\\
&\alpha_2&c_2&&&\\
&\ddots&\ddots&\ddots&&\\
&&\ddots&\ddots&\ddots&\\
&&&&\alpha_{n-1}&c_{n-1}\\
&&&&&\alpha_n
\end{bmatrix}
\]
By applying Doolittle's method, it can be shown that
\[
\alpha_1 = a_1
\]
and for $k = 2, 3,\ldots, n$,
\begin{align*}
\beta_k &=\frac{b_k}{\alpha_{k-1}}\\
\alpha_k&=a_k - \beta_k c_{k-1}
\end{align*}
\end{theorem}
\begin{theorem}[Thomas's Algorithm]
\hfill\\\normalfont If $\mathbf{Ax}=\mathbf{g}$ and define $\mathbf{Ux} = h$ where $\mathbf{x}:=(x_1,x_2,\ldots, x_n)^\text{T}$, $\mathbf{g}=(g_1,g_2,\ldots, g_n)^\text{T}$ and $\mathbf{h}=(h_1,h_2,\ldots, h_n)^\text{T}$. 
\begin{align*}
h_1&=g_1\\
h_i&=g_i-\beta_ih_{i-1},\;\;\; i = 2, 3, \ldots, n
\end{align*}
and
\begin{align*}
x_n &=\frac{h_n}{\alpha_n}\\
x_i &=\frac{h_i-c_ix_{x+1}}{\alpha_i}
\end{align*}
This method is known as \textbf{Thomas algorithm}.
\end{theorem}
\subsection{Pivoting Strategies}
\subsubsection{Partial Pivoting}
\begin{definition}[Partial Pivoting]\hfill\\\normalfont
At the $(k+1)$th step of the Gaussian Elimination process, $k = 0,1,\ldots, n-2$, choose the element having maximum absolute value in the $(k+1)$th column of $\mathbf{A}^{(k)}$ that lies on or below the diagonal so that
\[
|a_{s,k+1}^{(k)}|=\max_{i}|a_{i,k+1}^{(k)}|\;\;\;k+1\leq i\leq n
\]
and interchange row $k+1$ with row $s$.
\end{definition}
\subsubsection{Scaled Partial Pivoting}
\begin{definition}[Scaled Partial Pivoting]\hfill\\\normalfont
In the beginning, calculate $s_i$ for all row $i$, where
\[
s_i = \max_{1\leq j\leq n}\{|a_{ij}|\} , i = 1, \ldots, n
\]
At the $(k+1)$th step of the Gaussian Elimination process, $k = 0,1,\ldots, n-2$, choose $r$th row of $\mathbf{A}^{(k)}$ that lies on or below the diagonal where $r$ is determined by
\[
\frac{|a_{r,k+1}^{(k)}|}{s_r}=\max_{i}\frac{|a_{i,k+1}^{(k)}|}{s_i}\;\;\;k+1\leq i\leq n
\]
and exchange row $k+1$ with row $r$ and $s_{k+1}$ with $s_{r}$.
\end{definition}
\clearpage
\section{Error of Approximation}
\begin{definition}[Error]
\hfill\\\normalfont Given $n+1$ data values $(x_0,f_0), (x_1,f_1),\ldots, (x_n,f_n)$, we may define the error of approximation, $E$, by
\[
E=\sum_{i=0}^n|p(x_i)-f_i|
\]
If we take $p(x)$ to be a polynomial of degree $n$ then by choosing
\[
p(x_i)=f_i, i= 0,1,\ldots, n
\]
we can make $E = 0$. The function $p(x)$ is the \textbf{Lagrange interpolating polynomial}.\\
Alternatively, when $n$ is large, we may fit a polynomial of degree $k$ where $k\ll n$. We use
\[
E=\sum_{i=0}^n[p(x_i)-f_i]^2
\]
By minimising $E$ we achieve a \textbf{least square fit}.\\
If the function $f(x)$ is continuous specified on the interval $[a,b]$, we may define the error $E$ in approximating $f(x)$ by $p(x)$ as
\[
E=\max_{a\leq x \leq b}|p(x)-f(x)|
\]
It is common in practice to take $p(x)$ as a polynomial and minimise $E$ with respect to variations in the coefficients. This generate the \textbf{minimax} polynomial approximation.
\end{definition}
The following theorem justifies the choice of polynomial for approximation
\begin{theorem}[Weierstrass Approximation Theorem]
\hfill\\\normalfont Suppose $f$ is defined and continuous on $[a,b]$. For each $\varepsilon>0$ there exists a polynomial $P(x)$, defined on $[a,b]$, with property that
\[
|f(x)-P(x)|<\varepsilon\;\;\;\forall x\in[a,b]
\]
\end{theorem}
\subsection{Least Square Approximation}
\subsubsection{Discrete Data}
Given a set of $m+1$ discrete data points $(x_0,f(x_0)), (x_1,f(x_0)),\ldots (x_m,f(x_m))$, each carrying weight $w_0,w_1,\ldots, w_m$ respectively. \\
The task is to find an approximating polynomial function\footnote{The subscript $n$ denotes the power of the polynomial; thus the polynomial $p_n$ has $n+1$ unknwons}
\[
p_n(x;a_0,a_1,\ldots, a_n):=a_0+a_1x+\cdots+a_nx^n
\] 
$(m>n)$ such that
\[
E(a_0,a_1,\ldots, a_n)=\sum_{i=0}^m w_i[f(x_i)-p_n(x_i; a_0,a_1,\ldots, a_n)]^2
\]
is minimised with respect to the parameters $a_0, a_1,\ldots, a_n$.\footnote{$E$ can be understood as a weighted sum of squares of residue for each data.}\\
We require
\[
\frac{\partial E}{\partial a_j} = 9 \forall j = 0,1,\ldots, n
\]
Therefore, by differentiating with $a_j$,
\[
\sum_{i=0}^m w_i[f(x_i)-p_n(x_i; a_0,a_1,\ldots, a_n)]\frac{\partial}{\partial a_j}\left(-p_n(x_i;a_0,a_1,\ldots, a_n)\right)\;\;\;\forall j = 0,1,\ldots, n
\]
Rearranging, we have
\[
\sum_{i=0}^mw_i\underbrace{(a_0+a_1x_i+\cdots+a_nx_i^n)}_{p_n(x_i; a_0,a_1,\ldots, a_n)}x_i^j   =\sum_{i=0}^m w_if(x_i)x_i^j\;\;\;\forall j = 0,1,\ldots, n
\]
There is $n+1$ unknowns, and $n+1$ equations, so the linear system is\footnote{the $j$th column corresponding to the equation of specific $j$}
\[
\begin{bmatrix}
\sum_{i=0}^m w_ix_i^0&\sum_{i=0}^m w_ix_i^1&\cdots \sum_{i=0}^m w_ix_i^n\\
\sum_{i=0}^m w_ix_i^1&\sum_{i=0}^m w_ix_i^2&\cdots \sum_{i=0}^m w_ix_i^{n+1}\\
\vdots&\vdots&\ddots&\vdots\\
\sum_{i=0}^m w_ix_i^n&\sum_{i=0}^m w_ix_i^{n+1}&\cdots \sum_{i=0}^m w_ix_i^{2n}\\
\end{bmatrix}\begin{bmatrix}
a_0\\a_1\\\vdots\\a_n\end{bmatrix}=\begin{bmatrix}
\sum_{i=0}^m w_if(x_i)x_i^0\\\sum_{i=0}^m w_if(x_i)x_i^1\\\vdots\\\sum_{i=0}^m w_if(x_i)x_i^n\end{bmatrix}
\]
\subsubsection{Continuous Function}
Given a continuous function $f(x)$ defined on an interval $[a,b]$ with each value $x\in[a,b]$ associated with a weight $w(x)$. \\
The task is to find an approximating polynomial function $p_n(x;a_0,a_1,\ldots, a_n)$ such that
\[
E(a_0,a_1,\ldots, a_n)=\int_a^b w(x)[f(x)-p_n(x;a_0,a_1,\ldots, a_n)]^2\diff x
\]
is minimised with respect to the parameters $a_0,a_1,\ldots, a_n$.\\
The necessary condition sfor $E(a_0,a_1,\ldots, a_n)$ to be minimum are
\[
\frac{\partial E}{\partial a_j}=0 \forall j = 0,1,\ldots, n
\]
Differentiating with respect to $a_j$,
\[
\frac{\diff E}{\diff a_j} = -2\int_a^b w(x)[f(x)-p(x;a_0,a_1,\ldots, a_n)]\frac{\partial}{\partial a_j}(p_n(x;a_0,a_1,\ldots, a_n))\diff x = 0
\]
Rearranging,\footnote{$\frac{\partial}{\partial a_j}(p_n(x;a_0,a_1,\ldots, a_n))=x^j$ is obvious.}
\[
\int_a^b w(x)(a_0+a_1x+\cdots+a_nx^n)x^j\diff x = \int_a^b w(x)f(x)x^j\diff x\;\;\;\forall j = 0,1,\ldots, n
\]
\[
\sum_{k = 0}^n \int_a^b w(x)x^{k+j}\diff x = \int_a^b w(x)f(x)x^j\diff x\;\;\;\forall j = 0,1,\ldots, n
\]
Similarly, a matrix equation can be formed from this linear system of $n+1$ variables and $n+1$ equations. However, this matrix is \textbf{ill-conditioned}. 
\subsection{Lagrange Interpolation}
\begin{definition}\hfill\\\normalfont
Given a set of $n+1$ data points $(x_0,f_0),(x_1,f_1),\ldots,(x_n,f_n)$, we seek to find a polynomial $p_n(x)$ of degree $n$ which passes through each of the given points, i.e.
\[
p_n(x_i)=f_i\;\;\;i = 0,1,2,\ldots, n
\]
If we write $p_n(x)= a_0+a_1x+a_2x^2+\cdots+a_nx^n$, we have
\[
\begin{bmatrix}
1&x_0&\cdots&x_0^n\\
1&x_1&\cdots&x_1^n\\
\vdots&\vdots&\ddots&\vdots\\
1&x_n&\cdots&x_n^n
\end{bmatrix}\begin{bmatrix}a_0\\a_1\\\vdots\\a_n\end{bmatrix} = \begin{bmatrix}f_0\\f_1\\\vdots\\f_n\end{bmatrix}
\]
The determinant of the $(n+1)\times(n+1)$ Vandermonde matrix is $\displaystyle\prod_{0\leq i<j\leq n}(x_j-x_i)$, so the system has a unique solution as long as $x_i, i = 1,2,\ldots, n$ are all distinct.\\
The solved polynomial is the unique Lagrange interpolating polynomial.
\end{definition}
\begin{theorem}[Computation of Lagrange Interpolating Polynomial]
\hfill\\\normalfont Write
\begin{align*}
p_n(x)&=\alpha_0(x-x_1)(x-x_2)\cdots(x-x_n)\\
&+\alpha_1(x-x_0)(x-x_2)\cdots(x-x_n)\\
&+\cdots\\
&+\alpha_n(x-x_0)(x-x_1)\cdots(x-x_{n-1})
\end{align*}
Then, we have
\[
 \alpha_i = \frac{f_i}{\displaystyle\prod_{\substack{j=0\\ j\neq i}}^n(x_i-x_j)}
\]
So,
\[
p_n(x) = \sum_{i=0}^n l_i(x)f_i
\]
where
\[
l_i(x) = \prod_{\substack{j=0\\j\neq i}}^n\frac{x-x_j}{x_i-x_j}
\]
\end{theorem}
Note that $l_i(x)$ admits the property that
\begin{align*}
l_i(x_i)&=1\\
l_i(x_j)&\neq 1 \text{ for }j\neq i
\end{align*}
\begin{theorem}[$\psi(x)$ and alternative expression of Lagrange Interpolating Polynomial]
\hfill\\\normalfont Let $\psi(x)$ denote the polynomial
\[
\psi(x)=(x-x_0)(x-x_1)\cdots(x-x_n) = \prod_{i=0}^n (x-x_i)
\]
then, by differentiating on each product term once at a time, we have
\[
\psi^\prime(x) = \sum_{i=0}^n\prod_{\substack{j=0\\j\neq i}}^n (x-x_j)
\]
and hence
\[
\psi^\prime (x_i) = (x_i-x_0)\cdots(x_i-x_{i-1})(x_i-x_{i+1})\cdots(x_i-x_n)
\]
as all other terms containing $(x-x_i)$ will vanish.\\
Thus,\footnote{Do NOT forget multiply $(x-x_i)$ while using this formula!}
\[
l_i(x)=\frac{\psi(x)}{(x-x_i)\psi^\prime(x_i)}
\]
So $p_n(x)$ may be written as
\[
p_n(x)=\sum_{i=0}^n\frac{\psi(x)}{(x-x_i)\psi^\prime(x_i)} f_i
\]
\end{theorem}
\subsection{Truncation Error of Interpolating Polynomial}
\begin{theorem}[Extended Mean Value Theorem]
\hfill\\\normalfont Suppose that
\[
a\leq x_0<x_1<\cdots<x_k\leq b
\]
\[
f(x_0)=f(x_1)=\cdots = f(x_k)=0
\]
and $f(x),f^\prime(x),\ldots, f^{(k)}(x)$ are all continuous on $[a,b]$. There is a $\xi\in(x_0,x_n)\subset(a,b)$ such that
\[
f^k(\xi)=0
\]
\end{theorem}
\begin{theorem}[Lagrange Interpolating Polynomial Error Formula]
\hfill\\\normalfont Let $f(x)\in C^{(n+1)}[a,b]$. Further, let $p_n(x)$ interpolates $f(x)$ at $(n+1)$ distinct points $x+0,x_1,\ldots, x_n\in[a,b]$. Then, for $x\in[a,b], x\neq x_i, i = 0,1,\ldots, n$,
\[
f(x)-p_n(x)=\psi(x)\frac{f^{(n+1)}(\xi)}{(n+1)!}
\]
for some $xi\in\text{Spr}\{x,x_0,x_1,\ldots,x_n\}$.\\
Here, $\text{Spr}\{x,x_0,x_1,\ldots, x_n\}$ denotes the smallest interval containing $x,x_0,x_1,\ldots,x_n$.
\end{theorem}
It follows that an upper bound for the error is
\[
\delta(f(x)):=|f(x)-p_n(x)| \leq \frac{M}{(n+1)!}|\psi(x)|
\]
where
\[
M=\max_{a\leq \xi\leq b}|f^{(n+1)}(\xi)|
\]
\end{document}